% \numexpr is an *internal quantity*, not a run of digits -- in integers too.
%
% It is a complete number and it ends the scan. Expanding it to characters and
% then carrying on looking for more digits reads whatever follows:
% `\count0=\numexpr 1\relax\five` came out 15 instead of 1.
%
% That is never merely a wrong number. Every comparison expl3 makes is
%
%   \if_int_compare:w \__int_eval:w #1#2\__int_eval:w #3
%     \__int_eval_end: \__prg_TF_true:w \fi: \use_ii:nn
%
% so the tokens after the expression are the conditional's own branch
% machinery. Scanning past it expanded \__prg_TF_true:w, which swallowed the
% \fi:, and every *false* \int_compare:nNnTF in the kernel took its true
% branch -- `1 < 0` allocating a count register is what surfaced.
%
% Three copies of this rule: the gullet's dimension scanner had it (since
% \count5\p@), the gullet's integer scanner did not, and neither did the
% stomach's. The stomach now calls the gullet's expression scanner so they
% cannot drift again.
\input plain
\def\M#1{\message{<#1>}}
\def\five{5}
\count0=\numexpr 1\relax\five
\M{a=\the\count0}
\count1=\numexpr 2\relax 7
\M{b=\the\count1}
\dimen0=\dimexpr 1pt\relax\five
\M{c=\the\dimen0}
\count2=3\five
\M{d=\the\count2}
\M{e=\number\numexpr 1\relax\five}
\M{f=\ifnum\numexpr 1\relax<\numexpr 0\relax\five T\else F\fi}
\M{g=\ifnum\numexpr 1\relax>\numexpr 0\relax\five T\else F\fi}
\count3=\numexpr 2*3\relax\five
\M{h=\the\count3}
\M{i=\the\numexpr\numexpr 2\relax+1\relax}
\count4=\numexpr 4\relax
\M{j=\the\count4}
\M{k=\ifnum\numexpr 7\relax=7 T\else F\fi}
\bye
